% ============================================================
% Paragraphs: Real-Market SPY Calibration Results
% For use in the Experiments section of the thesis/paper
% ============================================================

%--- Market Data and Pre-processing ---%
To evaluate the proposed calibration framework under realistic market conditions, we employ SPY (SPDR S\&P~500 ETF Trust) call option data from two consecutive trading days: 2022-09-01, used for calibration, and 2022-09-02, used for out-of-sample validation. The descriptive statistics of the two option chains are summarised in Table~\ref{tab:spy_data}. On both days, the underlying spot price lies near \$394--396 and the options carry approximately 105--106 calendar days to expiry ($\tau \approx 0.29$ years). Mid-quote prices (average of bid and ask) are adopted as market prices. Implied volatilities and Vega values are subsequently recovered via Black--Scholes inversion using \texttt{py\_vollib\_vectorized}, and contracts for which the inversion fails (below-intrinsic or zero-Vega quotes) are discarded.

The risk-free rate is obtained from the Federal Reserve's daily par yield curve spanning 2020--2023. Par yields at the standard maturity nodes (1, 2, 3, 6, and 12 months) are first converted to continuously compounded rates via $r_c = \ln(1 + R_{\text{par}} \cdot t)\,/\,t$, and the term structure is then interpolated at the exact option maturity using a cubic spline. This procedure yields date-specific, maturity-matched risk-free rates of 3.027\% (2022-09-01) and 2.986\% (2022-09-02), reflecting the rising interest-rate environment of late 2022. Options are further filtered to retain only those with positive trading volume, days-to-expiry between 40 and 300, and log-moneyness $\ln(K/S_0) \in [-0.15,\, 0.20]$. After cleaning, 10 near-the-money call contracts are selected from each day for calibration and validation, respectively.

%--- Same-Day Calibration ---%
Table~\ref{tab:spy_sameday} reports the same-day calibration performance, where each method is calibrated on and evaluated against the 2022-09-02 option chain. The proposed method achieves an implied volatility Mean Relative Error (IV MRE) of \textbf{0.16\%}, compared to \textbf{0.20\%} for Zhang~(2025), representing a 20\% relative improvement. Both methods employ multi-start Adam optimisation with 10 restarts and 300 gradient steps per restart, making the computational cost comparable. The proposed method calibrates directly in implied volatility space with a Vega-weighted objective and enforces the Feller condition ($2\kappa\theta > \sigma^2$) as a soft penalty, yielding physically consistent parameter estimates. Zhang~(2025) optimises a normalised price MSE without the Feller constraint; in price space, its same-day MRE is 4.46\%.

%--- Cross-Day Calibration ---%
Table~\ref{tab:spy_crossday} reports the cross-day calibration experiment, in which the Heston parameters calibrated on the 2022-09-01 chain are applied without re-calibration to predict the 2022-09-02 option prices and implied volatilities. This scenario directly measures the temporal stability of the calibrated parameters under a one-day horizon. The proposed method achieves a cross-day IV MRE of \textbf{0.96\%}, compared to \textbf{1.59\%} for Zhang~(2025)---a reduction of 0.63 percentage points (pp). In terms of degradation from the same-day baseline, the proposed method incurs an IV MRE increase of only 0.80\,pp, substantially smaller than the 1.39\,pp degradation of Zhang~(2025), indicating superior parameter stability over consecutive trading days. In price space, the proposed method's cross-day QuantLib price MRE is 11.33\% versus 2.88\% for Zhang~(2025); however, this comparison is confounded by the different moneyness ranges of the two days' option chains---the 2022-09-02 chain selected by Zhang~(2025) lies deeper in-the-money ($\ln(K/S_0) \in [-0.128,\,-0.032]$) compared to the near-the-money range used by the proposed method ($\ln(K/S_0) \in [-0.032,\,-0.001]$), where deep in-the-money options carry higher prices and thus exhibit smaller relative pricing errors. Overall, the results demonstrate that the proposed IV-space calibration framework, augmented by Feller regularisation, provides both competitive same-day fitting accuracy and improved cross-day parameter stability relative to the price-space baseline.
